% Appendix F: the model answers to Paper B, with marks, from exam/paper_b_solutions.md.

% The answers' headings are labelled ansb:q:N.
\renewcommand{\answerprefix}{ansb}

\chapter{Model Answers to Paper B}
\label{ansb:ch:answers}

Marks are shown per part. Method carries them: a correct rule with an arithmetic slip in it is worth
more than a correct number with no working, and later parts consume earlier answers, so an error
should be followed through rather than penalized twice.

Code answers are marked on the algorithm, not on the syntax. A candidate who writes correct
pseudocode loses nothing. The questions are in Appendix~\ref{exb:ch:paper}.

% ==========================================================================================
\answersection{1}{The model, and the rate that destroys it}{10}

\answerpart{a}{2}
\textbf{Traditional programming} is given the \textbf{input} and the \textbf{rules} and produces the
\textbf{output}. The rules are written by hand.

\textbf{Machine learning} is given the \textbf{input} and the \textbf{output} and produces the
\textbf{rules}. It is handed pairs of examples and finds the relationship itself. \worth{1}

\textbf{Why the traditional arrangement fails on images.} The input varies without bound. Two
photographs of the same object share almost no pixel values, so there is no finite set of
conditions on pixel values that recognizes the object; a rule that covers one image fails on the
next. Symbolic AI could handle chess, where every position can be described exactly, and could not
handle vision, speech or translation, where the input cannot be enumerated. \worth{1}

\answerpart{b}{3}
$\yp = kx + m$, $\delta = \yref - \yp$, $\Delta e = \delta \cdot \LR$, then
$m \mathrel{+}= \Delta e$ and $k \mathrel{+}= \Delta e \cdot x$, with $\LR = 0.5$.

\textbf{Step 1} ($x = 3$, $\yref = 7$):
\[
  \yp = 0, \quad \delta = 7, \quad \Delta e = 3.5, \quad m = 3.5, \quad k = 0 + 3.5(3) = 10.5
\]
\textbf{Step 2} ($x = 4$, $\yref = 9$):
\begin{gather*}
  \yp = 10.5(4) + 3.5 = 45.5, \quad \delta = 9 - 45.5 = -36.5, \quad \Delta e = -18.25 \\
  m = 3.5 - 18.25 = -14.75, \quad k = 10.5 - 18.25(4) = -62.5
\end{gather*}
\textbf{Step 3} ($x = 3$, $\yref = 7$):
\begin{gather*}
  \yp = -62.5(3) - 14.75 = -202.25, \quad \delta = 7 + 202.25 = 209.25, \quad \Delta e = 104.625 \\
  m = -14.75 + 104.625 = \mathbf{89.875}, \quad k = -62.5 + 104.625(3) = \mathbf{251.375}
\end{gather*}
\credit{1 mark per step}

$\lvert\delta\rvert$ runs $7 \rightarrow 36.5 \rightarrow 209.25$: it grows by roughly a factor of
5 to 6 every step, and alternates in sign. This is not slow convergence, it is divergence. The
parameters are further from $(2, 1)$ after three steps than they were at $(0, 0)$.

\answerpart{c}{3}
Let $p = kx + m$ and $\delta = \yref - p$. One optimization step gives
\[
  m' = m + \delta\,\LR, \qquad k' = k + \delta\,\LR\,x
\]
so the new prediction for the \textbf{same} input is
\[
  p' = k'x + m' = (k + \delta\,\LR\,x)x + m + \delta\,\LR = p + \delta\,\LR\,(1 + x^2)
\]
and therefore
\[
  \delta' = \yref - p' = \delta - \delta\,\LR\,(1 + x^2)
    = \boxed{\;\delta\left[1 - \LR(1 + x^2)\right]\;}
\]
\worth{2}

The deviation shrinks when $\lvert 1 - \LR(1+x^2)\rvert < 1$, that is when
$0 < \LR(1+x^2) < 2$:
\[
  \LR < \frac{2}{1 + x^2}
\]
At $x = 4$: $\LR < 2/17 = \mathbf{0.1176}$. \worth{1}

Note what the formula says: the bias contributes the 1 and the weight contributes the $x^2$,
because the weight is adjusted by $\Delta e \cdot x$ and that adjustment is then multiplied by $x$
again on the next prediction. The loop gain grows with the \emph{square} of the input.

\answerpart{d}{2}
\textbf{The reasons the book states.} A learning rate of $0.0$ or less can never improve the model:
zero makes every adjustment zero, and a negative rate moves every parameter away from the reference
value. A learning rate of $1.0$ or more corrects by at least the full error on every step, which
makes training oscillate or diverge rather than converge. \worth{1}

\textbf{Why passing the check guarantees nothing.} The $\LR < 1.0$ bound coincides with the bound
from (c) only at $\lvert x \rvert = 1$, where $2/(1+x^2) = 1$. At $x = 0$, where the weight does not
move at all, the real bound is looser, $\LR < 2$; as soon as $\lvert x \rvert > 1$ it is tighter,
and it tightens quadratically. At $x = 4$ it is $0.1176$; the $\LR = 0.5$ used in (b) sails through
the check in \code{train()} and diverges anyway, which is exactly what (b) shows happening.

The check is necessary, not sufficient. What actually bounds the safe learning rate is the scale of
the input data, which \code{train()} never looks at. \worth{1}

% ==========================================================================================
\answersection{2}{Order, precision, and knowing when to stop}{10}

\answerpart{a}{2}
Trained in a fixed order every epoch, the model risks fitting the \textbf{ordering} rather than the
data: it sees the same sequence of corrections over and over and can settle into a pattern that
reflects how the training sets happen to be stored, rather than the relationship between $x$ and
$y$. Local minima are not the issue here: the book notes that a linear model has none, since its
squared-error surface has exactly one minimum. That benefit of shuffling arrives with the
non-linear networks from Chapter~\ref{c3:ch:nn} on, where it helps the parameters escape poor
regions. \worth{1}

What the shuffle protects is the model's generalization: reshuffling every epoch exposes the model
to the same data in varying order, so nothing in the storage order can be learned. Shuffling once
at construction would not do it; one fixed order is still one fixed order. \worth{1}

\answerpart{b}{3}
The \code{static} local (\code{initialized}) keeps its value between calls, so the body that seeds
the generator runs on the first call and never again. \worth{1}

Without it, every call would re-seed. \code{std::srand(std::time(nullptr))} has one-second
resolution, so several calls in the same second would reset the generator to the identical state
and \code{std::rand()} would return the identical sequence each time; the shuffle would stop
shuffling. \worth{1}

\textbf{What would have broken in the first arrangement.} An anonymous namespace gives every
translation unit its own private copy of the function \emph{and its own private \code{initialized}
flag}. With two \file{.cpp} files that each need the generator, the seeding happens once per file
rather than once per program: the second file's call re-seeds a generator the first file had already
set up, and the ``exactly once'' guarantee is gone. Seeding has to happen once per
\textbf{program}, which is what a single definition in \file{utils.cpp}, shared through
\file{utils.hpp}, gives. \worth{1}

\answerpart{c}{3}
\begin{gather*}
  \lvert 2 - 1.85\rvert = 0.15, \quad \lvert 4 - 3.9\rvert = 0.1, \quad
  \lvert 6 - 6.3\rvert = 0.3, \quad \lvert 8 - 8.1\rvert = 0.1 \\
  \MAE = \frac{0.15 + 0.1 + 0.3 + 0.1}{4} = \frac{0.65}{4} = 0.1625 \\
  \text{precision} = 1.0 - 0.1625 = \mathbf{0.8375}
\end{gather*}
\worth{1}

\textbf{Why the threshold must be strictly inside $(0.0, 1.0)$.}
\begin{itemize}
  \item \textbf{At $1.0$ or above:} unreachable. Precision is $1 - \MAE$ and the mean
    \emph{absolute} error cannot be negative, so precision can never exceed $1.0$ and reaches it
    only on an exactly perfect model. Training would always run the full epoch count and the early
    stop would never fire.
  \item \textbf{At $0.0$ or below:} it accepts a model whose mean absolute error is $1.0$ or worse,
    which defeats the point of checking.
\end{itemize}
\worth{1}

\textbf{A negative precision.} Predictions $3, 7, 1, 15$ against $2, 4, 6, 8$ give absolute errors
$1, 3, 5, 7$, so $\MAE = 4$ and precision $= 1 - 4 = \mathbf{-3}$.

What this reveals: \textbf{precision here is not a percentage and not a bounded score.} It is
$1 - \MAE$ measured in the \emph{units of the output data}, so it depends entirely on the scale of
$y$. On data whose outputs run into the thousands, even a very good model reports a large negative
``precision''. It is a stopping heuristic tuned to the book's small integer training data, not a
general accuracy metric, and the name flatters it. \worth{1}

\answerpart{d}{2}
\textbf{What it saves.} Computing the precision scans every training set and performs a prediction
for each. Over thousands of epochs that is a full extra forward pass over the whole data set per
epoch, which can be a large fraction of the total work when the training loop itself is cheap.
\worth{1}

\textbf{The two prices.}
\begin{enumerate}
  \item Up to nine epochs of unnecessary training after the threshold has actually been reached,
    since the check only fires on every tenth.
  \item The reported epoch count is quantized to the epochs the check runs after, 1, 11, 21 and so
    on, and so is not the true one. This is why the example in Chapter~\ref{c2:ch:training} prints
    ``reached after 11 epochs'' for data that in fact reached the threshold somewhere in epochs 2
    to 11.
\end{enumerate}
\worth{1}

% ==========================================================================================
\answersection{3}{Two output nodes}{14}

\answerpart{a}{3}
$x_1 = 1$, $x_2 = 0$.
\begin{align*}
  s_1 &= 0.2 + 1(0.3) + 0(0.6) = 0.5 &&\Rightarrow y_1 = \mathbf{0.5} \\
  s_2 &= 0.5 + 1(-0.4) + 0(0.2) = 0.1 &&\Rightarrow y_2 = \mathbf{0.1} \\
  s_3 &= -0.1 + 0.5(0.7) + 0.1(0.5) = -0.1 + 0.35 + 0.05 = 0.3 &&\Rightarrow y_3 = \mathbf{0.3} \\
  s_4 &= 0.4 + 0.5(-0.2) + 0.1(0.9) = 0.4 - 0.1 + 0.09 = 0.39 &&\Rightarrow y_4 = \mathbf{0.39}
\end{align*}
Every sum is positive, so ReLU passes all four through unchanged. \worth{3}

\answerpart{b}{4}
\textbf{Output layer:}
\begin{align*}
  \delta_3 &= Y_1 - y_3 = 0 - 0.3 = -0.3, & e_3 &= -0.3 \cdot 1 = \mathbf{-0.3} \\
  \delta_4 &= Y_2 - y_4 = 1 - 0.39 = 0.61, & e_4 &= 0.61 \cdot 1 = \mathbf{0.61}
\end{align*}
\worth{2}

\textbf{Hidden layer:} each hidden node feeds \textbf{both} output nodes, so its deviation is the
sum of two terms, one per output node, each weighted by the connection from that hidden node to it:
\begin{align*}
  \delta_1 &= e_3 w_5 + e_4 w_7 = -0.3(0.7) + 0.61(-0.2) = -0.21 - 0.122 = -0.332
    \Rightarrow e_1 = \mathbf{-0.332} \\
  \delta_2 &= e_3 w_6 + e_4 w_8 = -0.3(0.5) + 0.61(0.9) = -0.15 + 0.549 = 0.399
    \Rightarrow e_2 = \mathbf{0.399}
\end{align*}
\worth{2}

\textbf{How many terms, and why.} Two, one per node in the next layer. A hidden node's output was
sent to every output node, so it is partly responsible for every output node's error, and its share
along each path is that path's weight. In general $\delta_i = \sum_j e_j w_{ji}$ over all $j$ nodes
in the next layer. A candidate who used only one term has effectively disconnected one of the two
output nodes.

Note $\delta_2$: the two contributions have \textbf{opposite signs}, and the larger one wins.
Node~2 is being pulled down by output node~3 and up harder by output node~4.

\answerpart{c}{4}
\textbf{Output layer.} $\Delta c_3 = -0.3(0.1) = -0.03$, $\Delta c_4 = 0.61(0.1) = 0.061$. The
weights are scaled by the output layer's input, which is $y_1 = 0.5$ and $y_2 = 0.1$:
\begin{align*}
  b_3 &= -0.1 - 0.03 = \mathbf{-0.13}, & b_4 &= 0.4 + 0.061 = \mathbf{0.461} \\
  w_5 &= 0.7 - 0.03(0.5) = \mathbf{0.685}, & w_7 &= -0.2 + 0.061(0.5) = \mathbf{-0.1695} \\
  w_6 &= 0.5 - 0.03(0.1) = \mathbf{0.497}, & w_8 &= 0.9 + 0.061(0.1) = \mathbf{0.9061}
\end{align*}
\worth{2}

\textbf{Hidden layer.} $\Delta c_1 = -0.332(0.1) = -0.0332$, $\Delta c_2 = 0.399(0.1) = 0.0399$.
The weights are scaled by the network's input, $x_1 = 1$ and $x_2 = 0$:
\begin{align*}
  b_1 &= 0.2 - 0.0332 = \mathbf{0.1668}, & b_2 &= 0.5 + 0.0399 = \mathbf{0.5399} \\
  w_1 &= 0.3 - 0.0332(1) = \mathbf{0.2668}, & w_3 &= -0.4 + 0.0399(1) = \mathbf{-0.3601} \\
  w_2 &= 0.6 - 0.0332(0) = \mathbf{0.6}, & w_4 &= 0.2 + 0.0399(0) = \mathbf{0.2}
\end{align*}
\worth{2}

\answerpart{d}{2}
\begin{align*}
  s_1 &= 0.1668 + 0.2668 = 0.4336 \Rightarrow y_1 = 0.4336 \\
  s_2 &= 0.5399 - 0.3601 = 0.1798 \Rightarrow y_2 = 0.1798 \\
  s_3 &= -0.13 + 0.4336(0.685) + 0.1798(0.497) = -0.13 + 0.297016 + 0.089361
    = \mathbf{0.256377} \\
  s_4 &= 0.461 + 0.4336(-0.1695) + 0.1798(0.9061) = 0.461 - 0.073495 + 0.162917
    = \mathbf{0.550422} \\
  \delta_3 &= 0 - 0.256377 = \mathbf{-0.2564} \quad (\text{was } {-0.3}) \\
  \delta_4 &= 1 - 0.550422 = \mathbf{0.4496} \quad (\text{was } 0.61)
\end{align*}
Both deviations shrank in magnitude, so the step helped, and it helped node~4 substantially more
than node~3, which is what you would expect from the much larger error node~4 started with.
\worth{2}

\answerpart{e}{1}
$w_2$ and $w_4$ are unchanged, at $0.6$ and $0.2$. Both are the hidden layer's weights on input
$x_2$, and $x_2 = 0$ for this training set. The weight update is
$w \mathrel{+}= \Delta c \cdot \text{input}$, so a weight whose input is zero receives an adjustment
of exactly zero however large its node's error was.

This is not the dying ReLU problem. Nothing is broken: a weight that had no influence on this
prediction gets no share of the blame for it, which is correct. It will move on the next training
set that has $x_2 = 1$. \worth{1}

% ==========================================================================================
\answersection{4}{The network class, and what a stub can prove}{12}

\answerpart{a}{3}
\textbf{Two consequences of storing the layers by reference:}
\begin{enumerate}
  \item \textbf{\code{Shallow} does not own its layers.} The caller constructs them and must keep
    them alive for at least as long as the network; a network outliving its layers holds dangling
    references. It also means the class has nothing to destroy, which is why its destructor is
    \code{default}.
  \item \textbf{A layer can be replaced without touching \code{Shallow}.} This is precisely what
    Chapter~\ref{c5:ch:dense} does: \code{Stub} is swapped for \code{Dense} and not a line of the
    network class changes. It also lets a \emph{test} reach the layer after the network was built
    and manipulate it, which is what \code{setOutput()} exists for.
\end{enumerate}
Also creditable: reference members make the class non-assignable, which is consistent with the
book deleting the copy and move operators anyway. \credit{2 marks for any two}

\textbf{\code{predict()} returns \code{myOutputLayer.output()}}, a reference to the output layer's
own output vector. No member variable is needed because the output layer already stores it and the
network holds that layer by reference: the value is always live. Keeping a copy would introduce the
possibility of it going stale, which is the exact failure \code{setOutput()} is designed to expose.
\worth{1}

\answerpart{b}{3}
The three steps, per training set: \textbf{feedforward}, then \textbf{backpropagation}, then
\textbf{optimization}. \worth{1}

\textbf{Constraint 1: between the steps.} Backpropagation computes each node's error from the
outputs feedforward just produced, so it cannot run first. Optimization adjusts parameters using the
errors backpropagation just computed, so it cannot run before that. The three are a strict chain.
\worth{1}

\textbf{Constraint 2: inside backpropagation.} The output layer must go first. The hidden layer's
deviation is $\sum_j e_j w_{ji}$ over the output layer's \emph{computed errors}, which do not exist
until the output layer has compared its output against the reference value. \worth{1}

Worth noting what is \emph{not} constrained: the two \code{optimize()} calls can be made in either
order. Optimizing the hidden layer does not change \code{myHiddenLayer.output()}; only
\code{feedforward()} writes that, so the output layer still receives the values it actually
consumed.

\answerpart{c}{4}
\textbf{\code{setOutput()}} sets every element of the stub's output to a chosen value. It makes it
testable that \code{predict()} reads its output layer \textbf{live}: build a network, change the
output layer's output, and the network's prediction must change with it. Nothing on
\code{Interface} can establish this, because a network that cached a copy at construction would
satisfy every interface method identically; it would return a plausible vector of the right size,
and only an \emph{external} change to the layer can tell the two apart. \worth{1}

\textbf{\code{feedforwardCount()}} tallies calls to \code{feedforward()}. It makes the
\textbf{shape of the training loop} testable: \code{train()} must perform exactly one feedforward
per training set per epoch, so the count after training must be
$\mathit{sets} \times \mathit{epochs}$. Nothing else can establish it. A loop that runs a single
pass instead of every epoch lines up dimensionally, touches no invalid memory, and returns
\code{true} exactly as a correct loop does; without a tally the two are indistinguishable.
\worth{1}

Note that these two work because the stub \emph{computes nothing}. A layer whose output is a fixed
constant is what makes the network's wiring visible: any change in what the network predicts must
have come from the wiring, because it cannot have come from the arithmetic. \worth{1}

\textbf{Why the counter is incremented before the size check.} The tally then records how often the
layer was \textbf{asked} to feed forward, not how often it agreed to. That is the quantity a test of
the training loop actually wants: the loop's job is to make the right number of calls, and whether
each call was accepted is a separate question, tested separately by the return value. Incrementing
after the check would conflate the two, and a loop that made the right calls with wrongly sized data
would under-report rather than fail. \worth{1}

\answerpart{d}{2}
\textbf{The rule:} \code{std::terminate()} is reserved for the constructor. Everything else reports
failure through its return value. \worth{1}

\textbf{The reason:} a constructor has no return value, so it has no way to tell the caller that
construction failed. A \code{Shallow} built with no training data, or with an output layer whose
weight count does not match the hidden layer's node count, can never do anything useful, and
letting it exist only defers the failure to somewhere less informative. \code{train()} does have a
return value, so it can hand the decision back to the caller, who may reasonably want to retry,
with a different epoch count or threshold or simply again, rather than have the program killed;
\code{train()} resets both layers first, so every retry starts from freshly drawn parameters.
\worth{1}

% ==========================================================================================
\answersection{5}{The dense layer, and the mistake only Tanh can catch}{12}

\answerpart{a}{3}
\begin{booktable}{@{}p{10.6em}p{11.2em}L@{}}
  \thead{Member} & \thead{Size} & \thead{Quantity from Chapter~\ref{c3:ch:nn}} \\
  \midrule
  \code{myOutput} & \code{nodeCount()} & $y$, each node's output, after the activation \\
  \code{myPreActivationOutput} & \code{nodeCount()} & $s$, each node's weighted sum, before it \\
  \code{myError} & \code{nodeCount()} & $\Delta e$, each node's computed error \\
  \code{myBias} & \code{nodeCount()} & $b$, each node's bias \\
  \code{myWeights} & \code{nodeCount()} $\times$\newline \code{weightCount()} & $w$, each node's
    weights \\
  \code{myActFunc} & none & $\sigma$, the layer's activation function \\
\end{booktable}
\textit{(3 marks; deduct for a missing \code{myPreActivationOutput}, which is the one candidates
forget and the one part (c) is about.)}

\answerpart{b}{5}
\textbf{Feedforward.}
\begin{gather*}
  s = b + w_0 x_0 + w_1 x_1 = 0.3 + 0.6(1.0) + 0.6(1.0) = \mathbf{1.5} \\
  y = \tanh(1.5) = \mathbf{0.905148}
\end{gather*}
\worth{1}

\textbf{Backpropagation.}
\begin{gather*}
  \delta = \yref - y = 0.5 - 0.905148 = \mathbf{-0.405148} \\
  \sigma'(s) = 1 - \tanh^2(1.5) = 1 - 0.905148^2 = 1 - 0.819293 = \mathbf{0.180707} \\
  \Delta e = \delta \cdot \sigma'(s) = -0.405148 \times 0.180707 = \mathbf{-0.073213}
\end{gather*}
\worth{2}

\textbf{Optimization.} $\Delta e \cdot \LR = -0.073213 \times 0.2 = -0.014643$.
\begin{gather*}
  b = 0.3 - 0.014643 = \mathbf{0.285357} \\
  w_0 = 0.6 - 0.014643(1.0) = \mathbf{0.585357}, \qquad
  w_1 = 0.6 - 0.014643(1.0) = \mathbf{0.585357}
\end{gather*}
Both weights move by the same amount because both inputs are 1.0. \worth{1}

\textbf{Verification.}
\begin{gather*}
  s = 0.285357 + 0.585357 + 0.585357 = 1.456072 \Rightarrow y = \tanh(1.456072) = 0.896887 \\
  \delta = 0.5 - 0.896887 = \mathbf{-0.396887} \quad (\text{was } {-0.405148})
\end{gather*}
\worth{1}

The deviation moved by only $0.008$ on an error of $0.41$, about 2\%. This is tanh saturation doing
exactly what it is supposed to: at $s = 1.5$ the derivative is already down to $0.18$, so a large
error produces a small step. A candidate who observes that this node will take many epochs to move,
and that this is a feature rather than a fault, has understood the question.

\answerpart{c}{4}
Passing the \textbf{output} instead of the weighted sum:
\begin{gather*}
  \sigma'(y) = 1 - \tanh^2(0.905148) = 1 - 0.718795^2 = 1 - 0.516666 = \mathbf{0.483334} \\
  \Delta e_{\mathrm{wrong}} = -0.405148 \times 0.483334 = \mathbf{-0.195822} \\
  \text{ratio} = \frac{-0.195822}{-0.073213} = \mathbf{2.675}
\end{gather*}
\worth{2}

The step is 2.7 times too large. And the error is worst where it matters most: as
$\lvert s \rvert$ grows the true derivative collapses toward zero while the mistaken one flattens
out at a much larger value, so the bug removes precisely the braking that tanh saturation is there
to provide.

\textbf{Why ReLU hides it.} ReLU's derivative tests only the predicate ``is this value positive'',
and $\max(0, s)$ is positive exactly when $s$ is. So \code{actFuncDelta(Relu, s)} and
\code{actFuncDelta(Relu, y)} return the same number for every $s$, not merely for the values that
happen to be tested. There is no input that distinguishes them. \worth{1}

\textbf{The consequence for testing.} A suite that exercises only \code{Relu} cannot detect the bug
at all; every assertion passes and the layer is wrong. The test has to be written against
\code{Tanh} specifically, which is what \code{BackpropagateUsesPreActivationDerivative} in the test
suite of Chapter~\ref{c5:ch:dense} does. The general point: a test has to include a case that
\emph{distinguishes} the correct implementation from the plausible wrong one, not merely a case both
of them survive. \worth{1}

% ==========================================================================================
\answersection{6}{Padding, sharing, and the kernel's gradient}{14}

\answerpart{a}{3}
$\mathrm{pad} = 3 / 2 = 1$, so the padded input is $6 \times 6$:

\begin{console}
0 0 0 0 0 0
0 0 1 1 0 0
0 1 1 1 1 0
0 1 1 1 1 0
0 0 1 1 0 0
0 0 0 0 0 0
\end{console}
\worth{1}

With \code{ActFunc::None} and a bias of 0, the output is the raw weighted sum. The kernel's middle
column is all zeros, so each cell is (sum of the window's left column) minus (sum of its right
column):
\[
  s(0,0) = (0 + 0 + 0) - (0 + 1 + 1) = -2, \qquad s(0,3) = (0 + 1 + 1) - (0 + 0 + 0) = 2
\]
\textbf{Output:}

\begin{console}
-2  -1   1   2
-3  -1   1   3
-3  -1   1   3
-2  -1   1   2
\end{console}
\worth{1}

\textbf{What the kernel responds to.} It is a \textbf{vertical edge detector}. It compares the three
pixels down the left of the window against the three down the right and ignores the middle column
entirely, so it reacts to a left-right change in brightness and not at all to a top-bottom one.

\textbf{What the sign means.} Positive means brighter on the left of the window than on the right;
negative means the reverse. Zero means the neighbourhood is left-right symmetric, which is why the
output above is antisymmetric across the image: this input is a symmetric blob, so the left half
gives negative responses and the right half gives the mirrored positive ones. The magnitude is the
strength of the edge. \worth{1}

\answerpart{b}{3}
Both cases pad to \textbf{$6 \times 6$}: $\mathrm{pad} = 3/2 = 1$ and $\mathrm{pad} = 2/2 = 1$ are
the same number, so the padded input is $4 + 2(1) = 6$ either way. \worth{1}

\textbf{With the $3 \times 3$ kernel}, output cell $(i, j)$ reads padded rows $i \ldots i+2$. The
largest output index is 3, so the largest padded index read is $3 + 2 = 5$, the last row and
column. Everything is used.

\textbf{With the $2 \times 2$ kernel}, output cell $(i, j)$ reads padded rows $i \ldots i+1$. The
largest padded index read is $3 + 1 = 4$. \textbf{Padded row 5 and padded column 5 are never touched
by \code{feedforward()} at all}, and correspondingly never receive anything in
\code{backpropagate()}. \worth{1}

\textbf{Why, and what it implies.} A same-size output needs $\mathit{kernelSize} - 1$ extra rows in
total. For an odd kernel that is an even number and splits evenly, one on each side. For a
$2 \times 2$ kernel it is 1, but \code{pad = kernelSize / 2} adds 1 on \emph{each} side, which is
one too many; the loop starts at $(0,0)$ so the surplus lands at the bottom and the right and is
dead storage.

The alignment consequence is the real point. The padding is effectively one row and column
\emph{before} the image and none after, so output cell $(i,j)$ is computed from the window covering
input cells $(i-1, j-1)$ to $(i, j)$. There is no centre cell for an even kernel to sit on, so the
output is not centred on the input: a feature at input position $(r, c)$ reaches outputs $(r, c)$
to $(r+1, c+1)$, so it appears in the output shifted down and right by half a pixel. An odd kernel
has a true centre and does not suffer this. \worth{1}

\answerpart{c}{4}
With \code{ActFunc::None} the derivative is 1 everywhere, so \code{outGrad} equals the incoming
gradient at every position and no gradient is lost.
\[
  \text{biasGradient} = 2 + 4 + 6 + 8 = \mathbf{20}
\]
\worth{1}
\[
  \text{kernelGradient}[a][b] = \sum_{i,j} P[i+a][j+b] \cdot g[i][j]
\]
Only four gradients are non-zero: $g(0,1) = 2$, $g(1,3) = 4$, $g(2,0) = 6$, $g(3,2) = 8$. Each
kernel entry therefore has four terms. The centre entry, as a worked example:
\[
  k_{11} = P[1][2](2) + P[2][4](4) + P[3][1](6) + P[4][3](8) = 1(2) + 1(4) + 1(6) + 1(8) = 20
\]
\textbf{Kernel gradients:}

\begin{console}
12   14   14
12   20    8
 6    6    8
\end{console}
\worth{2}

\textbf{Updated parameters} at $\LR = 0.01$:
\[
  \text{bias} = 0.0 + 20(0.01) = \mathbf{0.2}
\]

\begin{console}
1.12   0.14  -0.86
1.12   0.20  -0.92
1.06   0.06  -0.92
\end{console}
\worth{1}

\answerpart{d}{4}
\textbf{Conv layer:} $3 \times 3 + 1 = \mathbf{10}$ trainable parameters.

\textbf{Dense layer} producing 16 outputs from 16 inputs: $16 \times 16 + 16 = \mathbf{272}$.
\worth{1}

\textbf{How many positions contributed to the centre kernel gradient.} \textbf{Four}, the four
output positions with a non-zero gradient, $(0,1)$, $(1,3)$, $(2,0)$ and $(3,2)$. In principle all
\textbf{sixteen} output positions contribute; twelve of them contributed a term of zero here only
because their gradients were zero. \worth{1}

\textbf{Why one update nevertheless.} The kernel is a \emph{single} set of nine numbers used at
every one of the sixteen output positions. Each position therefore gives a partial derivative with
respect to the same nine parameters, and the derivative of the total error with respect to one
parameter is the \textbf{sum} of those partials. Accumulating first and applying once is not an
optimization; it \emph{is} the gradient. \worth{1}

Applying an update per position instead would be a different algorithm and a wrong one: each of the
sixteen updates after the first would be computed against a kernel that the previous updates had
already moved, so the sum would no longer be the gradient at the point where it was evaluated. This
is the one place where a conv layer genuinely differs from a dense layer in its backward pass, and
it is a direct consequence of weight sharing: a parameter used in many places accumulates blame from
all of them. \worth{1}

% ==========================================================================================
\answersection{7}{A pooling layer that runs and returns true}{12}

\answerpart{a}{4}
\textbf{Fault 1: \code{>=} where \code{feedforward()} uses \code{>}.} \code{feedforward()} keeps a
value only when it is strictly greater than the best so far, so on a tie it records the
\textbf{first} occurrence. \code{backpropagate()} accepts equality, so it walks past the first
occurrence and records the \textbf{last}. On any block with a tied maximum the two halves disagree.

What the conv layer behind receives: a gradient at a position that did \textbf{not} produce the
pooled output, and nothing at the position that did. The gradient is real and lands on the wrong
pixel, so the conv layer updates its kernel from a neighbourhood the prediction never used.
\worth{2}

\textbf{Fault 2: \code{myInputGradients} is never reset to zero.} Each call writes exactly one cell
per block and leaves every other cell untouched, so those cells retain whatever the
\emph{previous} call left in them.

What the conv layer behind receives: this example's gradients mixed with stale gradients at up to
$\mathit{poolSize}^2 - 1$ positions per block, each left there by whichever earlier training example
last wrote that cell, possibly many examples ago. The layer never sees a clean gradient matrix after
the first training set. \worth{2}

\answerpart{b}{3}
\begin{booktable}{@{}p{6.4em}Lp{4.4em}p{4.4em}@{}}
  \thead{Block} & \thead{Values (position)} & \thead{First max} & \thead{Last max} \\
  \midrule
  top left & 0.5 (0,0), 0.5 (0,1), 0.2 (1,0), 0.3 (1,1) & (0, 0) & (0, 1) \\
  top right & $-0.1$ (0,2), $-0.4$ (0,3), $-0.6$ (1,2), $-0.9$ (1,3) & (0, 2) & (0, 2) \\
  bottom left & 1.0 (2,0), 0.4 (2,1), 0.1 (3,0), 0.9 (3,1) & (2, 0) & (2, 0) \\
  bottom right & 0.7 (2,2), 0.2 (2,3), 0.3 (3,2), 0.8 (3,3) & (3, 3) & (3, 3) \\
\end{booktable}

\noindent\textbf{What the code produces} (\code{myInputGradients} holding all zeros beforehand, as it
does on the first call after construction):

\begin{console}
 0   10   20    0
 0    0    0    0
30    0    0    0
 0    0    0   40
\end{console}

\noindent\textbf{What it should have produced:}

\begin{console}
10    0   20    0
 0    0    0    0
30    0    0    0
 0    0    0   40
\end{console}
\worth{2}

\textbf{They differ in the top-left block}, the only one with a tied maximum: the gradient 10 lands
on $(0,1)$ instead of $(0,0)$. Note how small the visible difference is, one value one column
across, and that the layer returns \code{true} either way. \worth{1}

\answerpart{c}{3}
\textbf{What is wrong.} \code{MaxPool} has \textbf{no trainable parameters}, and \code{myInput} is
not one. It is the last feedforward input, stored only so \code{backpropagate()} can re-locate which
cell held each block's maximum, and it is overwritten wholesale by the next \code{feedforward()}.
This method therefore adjusts something that is never learned and never read again: it burns
$O(n^2)$ work per training example to corrupt the layer's own record of what it was last shown.
\worth{1}

\textbf{The correct body:}

\begin{cppcode}
bool MaxPool::optimize(const double learningRate) noexcept
{
    (void)(learningRate);
    return true;
}
\end{cppcode}

\noindent\textbf{The general rule:} a layer with no trainable parameters has an \code{optimize()}
that is a \textbf{genuine no-op}, not a placeholder awaiting an implementation. There is nothing to
update, and inventing something to update is worse than doing nothing. \worth{1}

\textbf{How \code{flatten_layer::Interface} differs, and which is stronger.} It declares \textbf{no
\code{optimize()} at all}, not an empty one, not a defaulted one. The interface simply does not ask
for it.

That is the stronger design, because it moves ``this layer has nothing to optimize'' from a
convention into the type system: a \code{Flatten} that tried to optimize something would not
compile, whereas \code{MaxPool} relies on whoever writes it remembering that the method must do
nothing, exactly the memory lapse this listing shows. The book keeps \code{optimize()} on
\code{conv_layer::Interface} only because \code{Conv} and \code{MaxPool} share that interface and
\code{Conv} genuinely needs it; the no-op is the price of the shared interface, not a preference.
\worth{1}

\answerpart{d}{2}
\textbf{The symptom.} Almost nothing, on the surface. The network compiles, runs, predicts, and every
call returns \code{true}; there is no error message anywhere. Most training runs still end up
classifying all four digits, because the dense layer does most of the classifying and the conv
layer behind the pooling layer is only partly misled. What the user sees is a network that now and
then fails to learn: measured on the digits 0 to 3, a \code{MaxPool} with all three faults leaves
roughly one training run in seven unable to tell all four digits apart, where the correct layer
never fails. Retraining usually fixes it, which is exactly what makes the fault easy to ship.
\worth{1}

\textbf{The tests.} The test suite of Chapter~\ref{c10:ch:flatten} carries the \code{MaxPool} unit
tests of Chapter~\ref{c9:ch:maxpool} over unchanged, and they are what catches all three:
\begin{itemize}
  \item \textbf{Fault (a) 1, \code{>=}:} \code{BackpropagateRoutesToFirstMaxOnTies} fails, and so
    does \code{BackpropagateMatchesHandTrainedExample}, since the conv output from the
    hand-training exercise of Chapter~\ref{c6:ch:cnn} holds 1.9 three times in its top-right block.
  \item \textbf{Fault (c), the invented \code{optimize()}:} \code{OptimizeAcceptsAnyLearningRate}
    fails, because this \code{optimize()} validates a learning rate it has no use for and returns
    \code{false} for the zero, negative and out-of-range rates the test passes.
  \item \textbf{Fault (a) 2, the missing reset:} \code{BackpropagateResetsGradientsBetweenCalls}
    fails. It feeds two inputs whose maxima sit in different positions and backpropagates after
    each, so the gradients the first call wrote are still there after the second, where the correct
    layer has zeros. It is the only test that fails every time:
    \code{LearnsToRecognizeAllFourDigits}, the component test that trains on the digits 0 to 3 and
    checks the right output node wins for each, catches this fault only on some runs, roughly one
    in five. A test that backpropagated the same input twice would miss it entirely, since an
    assignment that writes the same cells twice looks clean.
\end{itemize}
\textit{(1 mark for a unit test that fails for each of the three faults.)}

% ==========================================================================================
\answersection{8}{Every layer, side by side}{8}

\answerpart{a}{4}
\begin{booktable}{@{}p{3.9em}LLLL@{}}
  \thead{Layer} & \thead{Trainable parameters} & \thead{\code{feedforward()}} &
    \thead{\code{backpropagate()}} & \thead{\code{optimize()}} \\
  \midrule
  \code{Fixed} & weight $k$, bias $m$ & $y = kx + m$ & $\delta = \yref - \yp$ &
    $k \mathrel{+}= \delta \cdot \LR \cdot x$,\newline $m \mathrel{+}= \delta \cdot \LR$ \\
  \addlinespace
  \code{Dense} & one bias and \code{weightCount()} weights per node &
    $y = \sigma\!\left(b + \sum w_i x_i\right)$ per node &
    $\Delta e = \delta \cdot \sigma'(s)$ per node, $\delta$ from the reference values or from the
    next layer &
    $b \mathrel{+}= \Delta e \cdot \LR$,\newline
    $w \mathrel{+}= \Delta e \cdot \LR \cdot \mathit{input}$ \\
  \addlinespace
  \code{Conv} & one shared kernel and one bias & as \code{Dense}, but local (one small window per
    output) and shared (the same kernel everywhere) & gradients accumulated across every position
    the kernel visited, plus input gradients for the layer behind &
    $\mathit{kernel} \mathrel{+}= \mathit{grad} \cdot \LR$,\newline
    $\mathit{bias} \mathrel{+}= \mathit{grad} \cdot \LR$ \\
  \addlinespace
  \code{MaxPool} & none & passes the largest value of each block forward & routes the incoming
    gradient to the position that held that maximum; every other position gets 0 &
    no-op: nothing to update \\
  \addlinespace
  \code{Flatten} & none & reshapes 2D to 1D, row-major & reshapes 1D back to 2D, the exact
    inverse & does not exist in the interface \\
\end{booktable}
\textit{(4 marks; roughly one per row beyond the first, with the \code{Fixed} row treated as free.)}

\answerpart{b}{2}
\textbf{Locality.} Each output looks only at a small window of the input rather than at all of it.

\emph{Buys:} the parameter count stops scaling with the input size, and the spatial structure of
the image is preserved rather than flattened away.
\emph{Costs:} a single layer cannot relate two distant parts of the image at all; reach has to be
built up by stacking layers, or by pooling between them to shrink the image under a fixed window.
\worth{1}

\textbf{Weight sharing.} The same kernel is reused at every position instead of each output owning
its own weights.

\emph{Buys:} the parameter count becomes independent of the image size entirely, and a pattern is
recognized wherever in the image it appears, which is what makes the layer generalize from far less
training data.
\emph{Costs:} the layer \emph{cannot} learn anything position-specific; it is forced to treat every
location identically. That is right for images and wrong for data where position itself carries
meaning. \worth{1}

\answerpart{c}{2}
\textbf{Feedforward}, left to right, each layer given the previous layer's \code{output()}:

\begin{console}
Conv::feedforward(input)
MaxPool::feedforward(conv.output())
Flatten::feedforward(maxPool.output())
Dense::feedforward(flatten.output())
\end{console}

\noindent\textbf{Backpropagation}, right to left. The target vector goes straight into the dense
layer; every earlier layer is given the next layer's \code{inputGradients()}:

\begin{console}
Dense::backpropagate(target)
Flatten::backpropagate(dense.inputGradients())
MaxPool::backpropagate(flatten.inputGradients())
Conv::backpropagate(maxPool.inputGradients())
\end{console}
\textit{(1 mark for both, including the target going into \code{Dense} and
\code{inputGradients()} carrying the rest.)}

\textbf{Optimization}, using gradients \code{backpropagate()} has already computed:

\begin{console}
Conv::optimize(learningRate)
MaxPool::optimize(learningRate)      // no-op
Dense::optimize(flatten.output(), learningRate)
Flatten                              // no optimize() exists
\end{console}

\noindent The dense layer needs the previous layer's \code{output()} as well as the learning rate,
because its weight update is scaled by the value that actually flowed along each weight. The order
of the \code{optimize()} calls does not matter: all the gradients already exist, and no layer's
optimization changes any other layer's stored output. \worth{1}

% ==========================================================================================
\answersection{9}{Off the laptop}{8}

\answerpart{a}{3}
Shapes: conv preserves $32 \times 32$ (it pads); pooling halves it to $16 \times 16$; flatten gives
$16 \times 16 = 256$ values; the dense layer has 64 nodes over those 256 inputs.

\begin{narrowtable}{@{}lr@{}}
  \thead{Layer} & \thead{Trainable parameters} \\
  \midrule
  conv & $3 \times 3 + 1 = 10$ \\
  max pooling & 0 \\
  flatten & 0 \\
  dense & $64 \times 256 + 64 = 16{,}448$ \\
  \midrule
  \textbf{total} & $\mathbf{16{,}458}$ \\
\end{narrowtable}
\worth{1}

\begin{narrowtable}{@{}lrr@{}}
  \thead{Type} & \thead{Bytes each} & \thead{Total footprint} \\
  \midrule
  \code{double} & 8 & 131,664 B (128.58 KB) \\
  \code{float} & 4 & 65,832 B (64.29 KB) \\
  \code{int8} & 1 & 16,458 B (16.07 KB) \\
\end{narrowtable}
\worth{1}

The budget is $64 \times 1024 = 65{,}536$ bytes.
\begin{align*}
  131{,}664 &> 65{,}536 \Rightarrow \textbf{double does not fit} \\
  65{,}832 &> 65{,}536 \Rightarrow \textbf{float does not fit either, by 296 bytes} \\
  16{,}458 &< 65{,}536 \Rightarrow \textbf{only int8 fits}
\end{align*}
The \code{float} case is the one that has to be done in bytes. ``About 64 KB against a 64 KB
budget'' reads like a fit and is a miss, by 296 bytes; that is before a single byte of program code,
stack, or input buffer. This is the argument for post-training quantization in one line: train at
full precision on a workstation, deploy at \code{int8}. \worth{1}

\answerpart{b}{3}
Q16.16 scales by $2^{16} = 65536$.
\begin{gather*}
  \text{\code{toFixed(0.75)}} = 0.75 \times 65536 = \mathbf{49{,}152} \\
  \text{\code{toFixed(-2.5)}} = -2.5 \times 65536 = \mathbf{-163{,}840}
\end{gather*}
\worth{1}
\[
  49{,}152 \times (-163{,}840) = \mathbf{-8{,}053{,}063{,}680}
\]
A 32-bit signed integer holds $-2{,}147{,}483{,}648$ to $2{,}147{,}483{,}647$. The product exceeds
that in magnitude by nearly a factor of four, so it must be formed in \code{std::int64_t} before
being scaled back down, which is precisely what the cast in \code{fixedMultiply()} is for. Two
values each scaled by $2^{16}$ multiply to a result scaled by $2^{32}$, so the operands can each be
small and the intermediate still overflow.
\[
  -8{,}053{,}063{,}680 \gg 16 = -122{,}880 \qquad
  \text{\code{toDouble}}(-122{,}880) = \frac{-122{,}880}{65{,}536} = \mathbf{-1.875}
\]
and $0.75 \times (-2.5) = -1.875$, exactly. \worth{1}
\[
  \text{resolution} = \frac{1}{65536} \approx \mathbf{0.0000153}
\]
with a representable range of roughly $\pm 32{,}768$. That is the trade: fixed dynamic range and
fixed absolute precision, in exchange for arithmetic that is a single integer instruction on a core
with no FPU rather than a software-emulated floating-point call. \worth{1}

\answerpart{c}{2}
\textbf{Two other constraints}, any two:
\begin{itemize}
  \item \textbf{Dynamic memory allocation.} Every layer in the book is built on \code{std::vector},
    which allocates on the heap. Long-running embedded systems risk heap fragmentation, allocation
    is not constant-time, and some bare-metal targets have no heap at all. The usual answer is
    \code{std::array} or plain arrays sized at compile time, trading runtime-sized layers for fully
    static memory use.
  \item \textbf{Real-time deadlines.} Inference inside a control loop has a hard budget; the book's
    example is a 1 kHz motor controller, giving sensor read, inference and actuation 1 ms in total.
  \item Also creditable: the need to quantize before deployment, and the absence of a filesystem or
    of \code{printf()} for diagnostics.
\end{itemize}
\worth{1}

\textbf{Multiply-accumulate count.} A conv layer performs roughly
$\mathit{inputSize}^2 \times \mathit{kernelSize}^2$ multiply-adds per feedforward pass:
\[
  32 \times 32 \times 3 \times 3 = \mathbf{9{,}216 \text{ MACs}}
\]
\textbf{Why worst case, not average.} A hard deadline is missed by a single late iteration; it is
not an average that has to be met but every instance. A network that is usually fast and
occasionally spikes still fails, and it fails in the way casual testing does not reveal, because the
average looks fine and the spike is rare. Scheduling has to be budgeted from the worst case, so the
worst case is the number that has to be measured. \worth{1}
